\section{Motivating Example}\label{sec:motivating-example}

\lai{Main point I want to make here is that clairvoyance semantics is
  impractical to compute but easier to translate into, while demand semantics is
  more practical to compute but hard to translate into directly. FLP lets you
  get the best of both worlds.}\ly{I'm not sure if that's true. We need actual
  examples of FLP compute these results in a reasonable time to support that.}

If it were simpler, one might want to translate between clairvoyance and demand semantics, particularly in the clairvoyance-to-demand direction. This is because translating a lazy function into clairvoyance semantics is relatively painless, while translating directly into demand semantics is much more challenging. \lai{say some more about why here, or give an example}.

For example, suppose we wanted to analyze the cost of following function in Curry, a lazy functional logic language which is syntactically similar to Haskell,\lai{Might need to explain this more...} 

[TODO: generic take n [1..] here, with a comment about what the [1..] syntax means]

Naively, one might assume that this function is non-terminating. With the advantage of laziness, though, only the first n elements of the list will be computed (as the rest are not returned as a part of the output). This time savings can be quantified using clairvoyance and demand semantics, and we expect this to result in a \(O(n)\) time cost.

Next, we can perform the transformation outlined
in~\cite{forking-paths} as follows, using the Tick monad to count function
calls:

[TODO: takeA here]

Notice how takeA is structurally similar to take. [TODO: expand on this more later].

Suppose now that we would prefer to analyze our program using demand semantics ~\cite{demand-semantics}\lai{Do we need to justify why here, or are we doing that earlier?}. This works in theory, but leaves us with some unresolved
issues in practice. The proof of equivalence to clairvoyance semantics in the paper does not include a simple and precise method for
translating our source function into a demand-semantics-based cost analysis
function. Indeed, this is one of the main drawbacks of the paper - the actual
translation is a brittle, manual process, which is time consuming and prone to
error if not mechanically proven to maintain functional correctness and time
cost equivalence. It is certainly possible, but untenable for most
applications.\lai{Do we need to do the work of showing insertD not defined in
  terms of insertA? I don't remember if we have that somewhere, and if not, that
  could be very annoying!}

Rather than defining insertD naively through transforming insertA, we can instead take full advantage of the nondeterminism \lai{Need to link this more clearly to the forward and backward thing from the intro} that Curry offers. Indeed, it turns out that the following is
equivalent to a directly translated version of insertD:

[TODO: Curry insertD goes here]

Now, all that remains is to prove that this version of insertD is also correct and
time cost equivalent to insert. How we do so is
detailed in sections 3-5\lai{TODO: Make this a link instead of hardcoded}. \lai{TODO: Too informal :/}

\begin{itemize}
  \item a more in-depth discussion of lazy evaluation semantics, esp.
        clairvoyance and demand semantics
  \item give an example that models each semantics in some programming language
        (probably Curry)
  \begin{itemize}
    \item this will require displaying/explaining the basic infrastructure
    \item what example? should be fairly small. take-append?
  \end{itemize}
  \item discuss how this is very annoying to write in demand semantics
  \item show how, in a logic language, you can ``write both in one program''
  \item this will need a discussion of what logic programming is, and probably a
        brief explanation of the specific language we are using (probably Curry)
\end{itemize}

% In the absence of recursion, call-by-need is semantically equivalent to
% call-by-name. However, the two semantics have different performance
% characteristics: in particular, call-by-need is \emph{cost-optimal} in the
% sense that a call-by-need derivation of \( t \Downarrow v \) will never be
% larger than a call-by-name derivation of the same, and may be smaller. Thus, a
% natural way to study call-by-name evaluation \emph{in paticular} is to augment
% the evaluation relation with an additional parameter that tracks the cost of
% evaluation.

% worked example
